% ============================================================
%  Expressões Algébricas — Apostila Premium | PratiqueJá
%  Engine: XeLaTeX
% ============================================================
\documentclass[12pt]{article}
\usepackage{../../pratiqueja}

% \hlm — destaque de resultado em modo-math (cor sem bold)
\newcommand{\hlm}[1]{{\color{darkblue}#1}}

\setsubject{Expressões Algébricas}

\begin{document}
\pagenumbering{arabic}
\setstretch{1.2}
\color{bodytext}

\heroheader{Expressões Algébricas}%
           {Termos, Valor Numérico e Operações}%
           {Matemática · Intermediário}

\setlength{\columnsep}{18pt}
\setlength{\columnseprule}{0.5pt}
\def\columnseprulecolor{\color{exborder}}

\begin{multicols}{2}

%─────────────────────────────────────────────────────────────
\TW{Def}{badgepurple}{Definição}{Variável, constante, termo e coeficiente}

\regra{Uma \textbf{expressão algébrica} combina números e letras
(variáveis) por operações. Cada parte separada por $+$ ou $-$ é um
\textbf{termo}.}

\exemp{%
  \textbf{Variável:} letra de valor desconhecido — $x, y, a$\\[2pt]
  \textbf{Constante:} número fixo — $3, -7, \tfrac{1}{2}$\\[2pt]
  \textbf{Coeficiente:} número que multiplica — em $5x^2$ é $5$\\[2pt]
  \textbf{Parte literal:} a variável com expoente — em $5x^2$ é $x^2$%
}

\exemp{%
  $3x + 2$ — binômio (2 termos)\\[2pt]
  $4x^2 - 3x + 7$ — trinômio (3 termos)\\[2pt]
  $-\dfrac{2}{5}ab^2$ — monômio (1 termo)%
}

%─────────────────────────────────────────────────────────────
\TW{V(x)}{darkblue}{Valor Numérico}{Substituir a variável e calcular}

\regra{O \textbf{valor numérico} é o resultado de substituir cada
variável por um número e calcular pela hierarquia das operações.}

\exemp{%
  \textbf{$E(x) = 2x^2 - 3x + 1$, \ $x = 2$:}\\
  $E(2) = 2\cdot2^2 - 3\cdot2 + 1$\\
  $= 8 - 6 + 1 = \hlm{3}$%
}

\exemp{%
  \textbf{$E(a,b) = a^2 + 2ab$, \ $a=3, b=-1$:}\\
  $= 3^2 + 2\cdot3\cdot(-1)$\\
  $= 9 - 6 = \hlm{3}$%
}

%─────────────────────────────────────────────────────────────
\TW{$\pm$}{darkblue}{Termos Semelhantes}{Mesma parte literal — podem ser somados}

\regra{Termos são \textbf{semelhantes} quando têm \emph{exatamente} a
mesma parte literal. Só semelhantes se somam/subtraem — basta operar
os coeficientes.}

\exemp{%
  \textbf{Simplificar:}\\
  $5x^2 - 3x + 2x^2 + 7x - 1$\\
  $= (5x^2 + 2x^2) + (-3x + 7x) - 1$\\
  $= \hlm{7x^2 + 4x - 1}$%
}

\exemp{%
  \textbf{Soma de polinômios:}\\
  $(3x^2 + 2x - 5) + (x^2 - 4x + 3)$\\
  $= \hlm{4x^2 - 2x - 2}$%
}

\nota{$3x^2$ e $3x$ \textbf{não} são semelhantes — partes literais
diferentes ($x^2 \neq x$).}

%─────────────────────────────────────────────────────────────
\TW{$\times$}{darkblue}{Multiplicação de Monômios}{Coeficientes $\times$ coef., literais $\times$ literais}

\regra{Multiplique os \textbf{coeficientes} e \textbf{some os
expoentes} das variáveis iguais.}

\exemp{%
  $3x^2 \times 4x^3 = \hlm{12x^5}$\\[3pt]
  $(-2ab^2)(5a^2b) = \hlm{-10a^3b^3}$\\[3pt]
  $2x \times 3y = \hlm{6xy}$ \;{\footnotesize(variáveis diferentes: justapostas)}%
}

%─────────────────────────────────────────────────────────────
\TW{Dist}{badgegreen}{Propriedade Distributiva}{Expandir e simplificar}

\formula{$a\cdot(b + c) = a\cdot b + a\cdot c$}

\small\color{bodytext}%
O fator externo multiplica \emph{cada} termo. Atenção ao sinal quando
ele é negativo.

\exemp{%
  $3x(2x - 5) = \hlm{6x^2 - 15x}$\\[3pt]
  \textbf{Fator negativo:}\\
  $-2(x^2 - 3x + 1) = \hlm{-2x^2 + 6x - 2}$\\[3pt]
  \textbf{Expandir e simplificar:}\\
  $2(x + 3) - 3(x - 1) = 2x + 6 - 3x + 3 = \hlm{-x + 9}$%
}

%─────────────────────────────────────────────────────────────
\TW{Aplic.}{badgegreen}{Exemplos Contextualizados}{Área e perímetro em linguagem algébrica}

\exemp{%
  \textbf{Área do retângulo:} largura $x$, comprimento $2x+3$.\\
  $A = x(2x + 3) = 2x^2 + 3x$\\
  Para $x = 4$: $A = 2(16) + 12 = \hlm{44}$%
}

\exemp{%
  \textbf{Perímetro:} triângulo de lados $a$, $a+2$, $2a-1$.\\
  $P = a + (a+2) + (2a-1) = \hlm{4a + 1}$%
}

\end{multicols}

\end{document}
